\documentclass[10pt,twocolumn]{article}

\usepackage[margin=0.75in,columnsep=0.25in]{geometry}
\usepackage{amsmath,amssymb,amsthm}
\usepackage{booktabs}
\usepackage{graphicx}
\usepackage[ruled,vlined,linesnumbered]{algorithm2e}
\usepackage{enumitem}
\usepackage[colorlinks=true,linkcolor=blue,citecolor=blue,urlcolor=blue]{hyperref}
\usepackage{microtype}

\newtheorem{theorem}{Theorem}
\newtheorem{proposition}{Proposition}
\newtheorem{corollary}{Corollary}
\newtheorem{lemma}{Lemma}
\newtheorem{definition}{Definition}
\newtheorem{assumption}{Assumption}
\newtheorem{remark}{Remark}
\newtheorem{problem}{Problem}
\newtheorem{example}{Example}

\newcommand{\Ex}{\mathbb{E}}
\newcommand{\hbm}{\mathsf{HBM}}
\newcommand{\cpu}{\mathsf{CPU}}
\newcommand{\nvme}{\mathsf{NVMe}}
\newcommand{\none}{\varnothing}
\newcommand{\Tcyc}{T_{\mathrm{cyc}}}
\newcommand{\Tgpu}{T_{\mathrm{gpu}}}
\newcommand{\Thost}{T_{\mathrm{host}}}
\newcommand{\Tser}{T_{\mathrm{ser}}}
\newcommand{\Txfer}{T_{\mathrm{xfer}}}
\newcommand{\Bmax}{B_{\max}}
\newcommand{\shat}{\hat{s}}
\newcommand{\eps}{\varepsilon}

\title{\textbf{Cost-Based Scheduling of Zero-Decode Semantic Filter Pipelines}}
\author{
  Anonymous Author(s)\\
  \textit{Draft of \today. Step-cost constants are calibrated on Qwen3 4B FP8 / H100 (measured companion: plans/cost\_model.md); the strategy-comparison numbers in Section 2.5 remain analytical.}
}
\date{}

\begin{document}
\maketitle

\begin{abstract}
Semantic filter pipelines apply a sequence of LLM-evaluated predicates to every document in a corpus. When each predicate is zero-decode, meaning the verdict is read from the logits of one forward pass with no generation, pipeline cost is dominated by prefill compute and by the movement of KV cache state. Serving engines execute such workloads as streams of independent requests, which ignores three sources of structure: each document's KV is reusable across filters, filters have selectivities that shrink downstream work, and the corpus is available upfront. We formalize the resulting joint batching and placement problem over one or more GPUs. Although the problem contains classical NP-hard scheduling and caching problems, in a first-order cost model a simple depth-first cohorting policy is provably within a factor $1 + z\eps$ of optimal, where $z$ is the pipeline depth and $\eps \approx 6 \times 10^{-3}$ is a measured hardware ratio, while the engine default loses up to $9\times$. The interesting scheduling problem therefore lives at the boundary of this guarantee, and we characterize that boundary exactly: streaming arrival, contested residency, documents exceeding single-device memory, warm KV persisted from earlier queries, and second-order step cost effects. For these regimes we develop a calibrated step cost model with provable identifiability requirements on its calibration design, a rolling-horizon program over aggregated document groups whose solve time hides behind GPU execution, and a length-aware placement policy derived from a closed-form crossover between transfer and recomputation. The device configuration, replica count and sharding degree, is an input throughout: replicas reduce to balanced independent instances inside the scheduler, and sharding enters only through the cost model's constants, leaving the cohorting guarantee unchanged.
\end{abstract}

\section{Introduction}
\label{sec:intro}

LLMs are increasingly integrated as operators in data systems, letting users express filters and maps over unstructured text in natural language, both through SQL extensions and through declarative pipeline frameworks~\cite{docetl,lotus,palimpzest}. The workhorse operation is the semantic filter: a natural-language predicate evaluated by an LLM on every document, with pipelines applying several filters in sequence. Consider the following example, which we carry through the paper.

\begin{example}[Contract screening]
\label{ex:contracts}
A compliance team screens a collection of 200{,}000 contracts with three questions applied in order: does the contract concern personal data; does it impose audit obligations; does it survive a change of control. Only contracts passing all three are kept. Each question is evaluated by placing the contract text before a short instruction suffix and reading the verdict from the model's output logits at the final position, so no text is generated.
\end{example}

Filters of this form are \emph{zero-decode}: an evaluation costs one forward pass, so at corpus scale the pipeline's cost is dominated by prefill compute and by the movement of KV cache state. A rich line of work reduces how many LLM calls such a pipeline makes, through model cascades with statistical guarantees~\cite{supg,bargain}, task cascades that also vary the operation and document fraction~\cite{taskcascades}, model selection~\cite{frugalgpt}, and classical predicate optimization~\cite{predmig,probpred}. This paper asks the question that remains after the call set is fixed: given the evaluations that must run, how should they be executed on the hardware so the pipeline finishes as fast as possible?

\begin{table}[t]
\centering
\small
\begin{tabular}{lccc}
\toprule
Approach & Reuse & Placement & Selectivity \\
\midrule
Engine default~\cite{vllm} & recency & reactive & unused \\
Prefix caching~\cite{sglang} & recency & reactive & unused \\
Offload systems~\cite{flexgen} & fixed order & planned & unused \\
FPS (this paper) & planned & planned & estimated \& used \\
\bottomrule
\end{tabular}
\caption{How execution approaches treat the three structures of a filter pipeline: KV reuse across filters, KV placement across tiers, and selectivity. The system estimates selectivity from a pilot sample and from its own verdicts, observations the other approaches receive and discard; Section~\ref{sec:selest} states where the estimates matter.}
\label{tab:approaches}
\end{table}

\textbf{Prior Work and Limitations.} Once the calls are fixed, execution falls to a serving engine, and Table~\ref{tab:approaches} summarizes how existing execution approaches treat the three structures a filter pipeline offers. Engines such as vLLM, Orca, and Sarathi-Serve~\cite{vllm,orca,sarathi} schedule streams of independent requests: KV reuse happens only when recency-based caching happens to retain a prefix, placement is reactive eviction under memory pressure, and selectivity plays no role (first row of Table~\ref{tab:approaches}). Prefix caching systems such as RadixAttention~\cite{sglang} retain prompt prefixes deliberately, but retention is still decided by recency rather than by whether a pipeline will return for the survivors (second row). Offload systems such as FlexGen~\cite{flexgen} plan placement, but for a fixed request order and without pipeline structure (third row). None of the three uses the filters' selectivities, even though selectivity is not privileged information: every verdict the pipeline produces is an observation of it, and these engines receive those observations and discard them. And none exploits the fact that the entire corpus sits in a queue that could be planned over. The cost of this mismatch is large: in Example~\ref{ex:contracts}, the KV for a 4K-token contract occupies 288\,MiB at FP8 precision, one H100 holds only a few hundred such contracts, and our analysis shows the engine default leaves up to $9\times$ of makespan on the table (Section~\ref{sec:motivation}).

\textbf{Plan the Execution, Not Just the Calls.} Our key insight is that a filter pipeline gives the execution layer structure that a stream of independent requests does not have, and that this structure can be planned over. In Example~\ref{ex:contracts}, the contract text is a shared prompt prefix for all three questions, so the KV cache built to answer the first question can be reused for the second and third on every contract that survives. The filters also have selectivities: if only 30\% of contracts concern personal data, then 70\% of the corpus never needs the later filters, and these rates can be estimated from a small sample. Finally, the whole corpus is available before execution starts, so the system can decide upfront which documents to admit, which to keep resident, and which to spill, instead of reacting to memory pressure.

\textbf{Challenges.} However, planning this execution is challenging for several reasons. First, two resources interact. Each engine step has a token budget that limits how much work fits in it, and the device has a memory budget that limits which documents can be evaluated at all, because evaluation requires the document's KV to be resident. A step can have token room but no evaluable documents, or the reverse, and every placement decision constrains the compositions available to all later steps. Second, the planning problem is computationally hard. Even simplified versions are NP-hard: choosing what runs in each step is a form of bin packing, and choosing where KV lives is a form of caching with variable-size pages, both classical NP-hard problems~\cite{uzsoy,cachinghard}. Exact per-document planning is also out of reach at corpus scale, since the number of decisions grows with the number of documents. Third, the selectivities are not known upfront, and they vary with document length, so the plan must work with estimates and improve them as verdicts arrive.

\textbf{Our Approach.} We first ask how large the differences between execution strategies can be, taking the classical execution strategies of query processing as candidates (Section~\ref{sec:motivation}). We make two observations. The engine default can be up to $9\times$ slower than the best strategy. And a simple strategy, depth-first cohorting, which is block-pipelined execution with the block sized by KV memory, is provably within a factor $1 + z\eps$ of any feasible schedule when the corpus is cold and step costs are first-order, where $z$ is the pipeline depth and $\eps \approx 6 \times 10^{-3}$ is the measured ratio of cached-read time to prefill time per token (Qwen3 4B FP8 on H100, at the 32-token suffix reference). So planned execution matters, a simple policy captures most of it, and that policy remains available to our scheduler, since it is a feasible point of the scheduler's program. The rest of our machinery targets the settings where the guarantee does not apply: the corpus streams in, another tenant shares the device, single documents do not fit in memory, KV from earlier queries is already in storage, or measured step costs deviate from the first-order model. We address the challenges as follows:
\begin{itemize}[leftmargin=*,itemsep=1pt]
\item \textbf{Joint step and placement planning.} Because the two budgets are coupled, the scheduler chooses each step's composition and each document's KV placement together, using a rolling-horizon program re-solved every step and overlapped with GPU execution so it adds no runtime. A calibrated cost model scores candidate steps, and we prove which of its terms calibration can identify; a closed-form crossover between transfer and recomputation makes placement length-aware and prices peer-device memory (Sections~\ref{sec:cost}--\ref{sec:placement}).
\item \textbf{Histogram-style grouping.} Per-document planning does not scale, so we summarize the corpus the way an optimizer summarizes a column: documents are bucketed by length, a group pairs a bucket with a pipeline position, and the scheduler plans counts over groups. We prove the relaxation we solve loses at most one document per group per step to rounding, a vanishing fraction of total work at corpus scale (Sections~\ref{sec:hardness}, \ref{sec:sched}).
\item \textbf{Selectivity estimation.} Per-group pass rates are estimated from a stratified pilot sample and updated online from the pipeline's own verdicts. We state where the estimates matter, replica balance, retention decisions, and the boundary regimes, and where they do not (Section~\ref{sec:selest}).
\end{itemize}

\textbf{Contributions.} To summarize, our contributions are as follows.
\begin{enumerate}[leftmargin=*,itemsep=1pt]
\item We formalize filter pipeline scheduling (FPS), parametric in the device configuration, and identify the structure that makes real instances tractable: grouping documents by filter and length bucket collapses the decision space to counts, with a provably small rounding loss (Sections~\ref{sec:problem}, \ref{sec:hardness}, \ref{sec:sched}).
\item We prove depth-first cohorting is $(1+z\eps)$-optimal in the first-order cost model, show the guarantee is invariant under tensor sharding, and characterize its envelope exactly, alongside an analytical $9\times$ gap for the engine default (Section~\ref{sec:motivation}).
\item We develop the step cost model with its identifiability theorem and feasibility-bounded calibration design (Section~\ref{sec:cost}).
\item We develop the replica-balanced rolling-horizon scheduler, with an overlap condition under which planning adds zero runtime, and the length-aware placement policy with its closed-form crossover (Sections~\ref{sec:sched}, \ref{sec:placement}).
\item We lay out an evaluation testing both sides of the envelope boundary on real corpora (Section~\ref{sec:eval}).
\end{enumerate}

We present the problem, the motivating analysis, and an overview of our approach in Section~\ref{sec:problem}. Section~\ref{sec:hardness} establishes hardness. Sections~\ref{sec:cost} through~\ref{sec:placement} develop the three components, Section~\ref{sec:eval} presents the evaluation plan, and Section~\ref{sec:related} covers related work. Proofs are deferred to Appendix~\ref{app:proofs}.

\section{Problem Definition}
\label{sec:problem}

\subsection{Setup}
\label{sec:setup}

\paragraph{Workload.} A corpus $D$ of $n$ documents is processed by a pipeline of $z$ semantic filters $F_1, \dots, F_z$ applied in order: a document is evaluated by $F_{o+1}$ only if it passes $F_o$. Documents are partitioned into $L$ length buckets with representative token lengths $h_l$. Filter $F_o$ is an instruction suffix of $c_o$ tokens, with $c_o$ small (16 to 64 in our workloads), and the prompt for evaluating $F_o$ on document $d$ is the tokens of $d$ followed by the suffix, so the document is a shared prompt prefix across filters. A filter is \emph{zero-decode} if its verdict is a function of the output logits at the final position of one forward pass; no tokens are generated. The selectivity $s_{o,l} \in [0,1]$ is the probability that a bucket-$l$ document at filter $o$ passes. Selectivities are estimated, either from a stratified pilot sample before planning or online during execution (Section~\ref{sec:selest}).

\paragraph{Running example.} We carry Example~\ref{ex:contracts} through the paper with concrete numbers: $n = 200{,}000$ contracts in buckets from 512 to 16K tokens, $z = 3$ filters of 32 tokens each, selectivities near $(0.3, 0.4, 0.5)$, served by Qwen3 4B at FP8 on H100s. One cached token costs $m = 72$\,KiB, so a 4K-token contract holds 288\,MiB of KV and free device memory of 68\,GiB holds about 240 such contracts.

\paragraph{Hardware and memory.} Each device serves the model with per-token KV cost $m$ bytes. KV lives in a tier set $\mathcal{K}$ with capacities $M_k$ and bandwidths $\beta_k$, drawn from $\{\hbm, \cpu, \nvme\}$ plus peer-device memory when replicas exist; $\none$ denotes discarded state, restorable only by recomputation. Evaluating a filter on a document requires that document's KV resident in $\hbm$.

\paragraph{Execution.} The engine runs synchronous steps. A step is a multiset $\sigma = \{(C_i, H_i)\}_{i=1}^{N}$ of requests, where request $i$ processes $C_i$ new tokens against $H_i$ cached tokens: a filter evaluation contributes $(c_o, h_l)$, and prefill work contributes $(C_i, H_i)$ with $C_i$ document tokens scheduled this step against $H_i$ already processed; prefill is divisible across steps while suffixes are atomic. Define the token count and attention pair count
\begin{equation}
\label{eq:BP}
B(\sigma) = \sum_i C_i, \qquad
P(\sigma) = \sum_i \Big( C_i H_i + \tfrac{C_i(C_i+1)}{2} \Big),
\end{equation}
with $B \le \Bmax$, the engine budget. Filter suffixes are atomic: a suffix runs within one step, since the verdict is read at its final position and suffixes are far below the token budget; only prefill is chunked. The cycle time of a step is
\begin{equation}
\label{eq:cycle}
\Tcyc(\sigma) = \max\big( \Tgpu(\sigma),\ \Thost(\sigma) + \Txfer(\sigma) \big) + \Tser,
\end{equation}
where $\Thost$ is host preparation, $\Txfer$ the non-overlapped part of tier transfers, and $\Tser$ a constant for logit readout and dispatch. Table~\ref{tab:notation} summarizes notation.

\begin{table}[t]
\centering
\small
\begin{tabular}{ll}
\toprule
Symbol & Description \\
\midrule
$D$, $n$, $z$, $L$ & corpus, its size, filter count, bucket count \\
$c_o$, $h_l$ & suffix length of $F_o$; token length of bucket $l$ \\
$s_{o,l}$, $\shat_{o,l}$ & true and estimated selectivity \\
$m$ & KV bytes per token \\
$\mathcal{K}$, $M_k$, $\beta_k$ & tier set; capacity and bandwidth of tier $k$ \\
$G$, $\nu$ & device count; sharding degree ($G/\nu$ replicas) \\
$\sigma$, $N$, $B$, $P$ & step; its request, token, and pair counts \\
$\Bmax$, $W$ & token budget; HBM workspace reserve \\
$g{=}(o,l)$, $n_{g,t}$, $m_{g,k,t}$ & group; its count and tier occupancy \\
$x_{g,u}$, $q_{l,u}$, $w_{g,k\to k',u}$ & evaluations, prefill tokens, transfers \\
$\tau$, $\Phi$ & horizon length; terminal cost \\
$t_{\mathrm{pre}}$, $t_{\mathrm{read}}$, $\eps$ & per-token prefill and cached-read \\
 & \quad times (measured); their ratio \\
$\Tcyc$ & cycle time of a step \\
\bottomrule
\end{tabular}
\caption{Notation.}
\label{tab:notation}
\end{table}

\subsection{Device Configuration}
\label{sec:devices}

The deployment is an input: $G$ identical accelerators arranged as $G/\nu$ data-parallel replicas of tensor-sharding degree $\nu$, fixed before scheduling begins. The configuration is consumed in three places downstream and nowhere else. First, replicas partition the corpus: each replica runs its own instance of the scheduler and the pipeline finishes when the slowest replica finishes, so the scheduler's first act is a balance step, exact up to one document per group (Section~\ref{sec:sched}). Second, sharding changes the cost model's constants and nothing structural: per-device KV per token becomes $m/\nu$, per-device compute per token becomes $2P_{\Theta}/(\nu\,\mathrm{FLOPS})$ up to all-reduce overhead, and free KV capacity aggregates across the group, so calibration simply runs on the deployed configuration (Section~\ref{sec:cost}). The mode is a genuine choice for any model that fits on one device, which includes Qwen3 32B at FP8 on an H100, and the trade-off is capacity against communication: replicas duplicate weights, so $G$ replicas hold $G(M_{\hbm} - W_{\Theta})$ of KV, while a degree-$\nu$ group stores weights once and holds $\nu M_{\hbm} - W_{\Theta}$, at the cost of all-reduce time per layer. For 32B on four H100s this is roughly 190\,GB of KV as replicas against 285\,GB as one sharded group. Third, when $G/\nu > 1$, a peer replica's free HBM is one more placement tier, reachable at NVLink bandwidth (Section~\ref{sec:placement}). Everything else in the paper is stated per replica.

\subsection{The Scheduling Problem}

A schedule is a sequence of steps with placement actions between them: transfers between tiers and evictions to $\none$. A schedule is feasible if every step respects its token budget, evaluates only $\hbm$-resident documents, respects filter precedence, and never exceeds any tier's capacity, counting KV written during the step. FPS asks for a feasible schedule processing every document to completion that minimizes expected makespan, the sum of cycle times, in expectation over pass outcomes; with $G > 1$ data-parallel devices, the maximum over devices of that sum. When selectivities are estimated online the problem is one of choosing a policy, and our algorithm is a certainty-equivalent policy for it.

\paragraph{Why this is a planning problem.} Two budgets interact. The token budget bounds work per step; the residency budget bounds which documents are evaluable at all. A step may have token headroom but no evaluable documents because their KV was evicted, or full residency and no headroom. Residency decisions constrain the compositions available to every later step, so good execution requires looking ahead, and the structure available for looking ahead (selectivities, the queue of pending documents, the length distribution) is exactly what a request-stream engine does not see.

\subsection{Motivating Analysis}
\label{sec:motivation}

We compare execution strategies under first-order cost arithmetic: prefill compute at the device's FLOP rate, KV reads at HBM bandwidth, transfers at tier bandwidths, and a per-step overhead. The candidate strategies are not new; they are the classical execution strategies of query processing, applied to KV state. Running the pipeline filter-at-a-time is operator-at-a-time execution~\cite{graefe-survey}, with the survivors' KV as the intermediate result, which can be discarded and rebuilt (\emph{breadth-discard}) or materialized in memory (\emph{breadth-retain}). Running it document-batch-at-a-time is block-pipelined execution in the MonetDB/X100 sense~\cite{monetx100}, with the block size set by KV memory rather than cache lines: this is \emph{depth-first cohorting}, which admits a cohort sized to fill free memory, runs all $z$ filters on it while its KV is hot, then admits the next. Materializing intermediates to slower storage (\emph{spill-return}) and fusing consecutive operators into one pass (\emph{fused})~\cite{hyper} complete the space. Table~\ref{tab:motivation} compares the \emph{engine default}, continuous batching with LRU prefix caching, against the two strategies that bracket this space, breadth-discard and cohorting, across five regimes spanning document length, selectivity, corpus size, and model scale. A wider analytical sweep over both model scales, lengths to 128K tokens, corpora to $10^7$ documents, and pipelines of depth ten preserves the ordering, with the default losing up to $9.2\times$.

\begin{table}[t]
\centering
\small
\begin{tabular}{l@{\hspace{6pt}}ccccc}
\toprule
 & \multicolumn{5}{c}{Slowdown vs.\ cohorting ($1.00$ = best)} \\
Strategy & R1 & R2 & R3 & R4 & R5 \\
\midrule
Depth-first cohorting & \textbf{1.00} & \textbf{1.00} & \textbf{1.00} & \textbf{1.00} & \textbf{1.00} \\
Breadth-discard       & 1.05 & 1.75 & 2.71 & 1.75 & 2.94 \\
Engine default        & 1.21 & 1.05 & 3.12 & 2.01 & 3.38 \\
\bottomrule
\end{tabular}
\caption{Analytical makespan as a slowdown factor (lower is better). R1: $h{=}512$, $s{=}0.05$, $n{=}10^6$, 4B. R2: $h{=}512$, $s{=}0.5$, $n{=}1800$, 4B. R3: $h{=}1$K, $s{=}0.9$, $n{=}10^5$, 4B. R4: $h{=}8$K, $s{=}0.5$, $n{=}2{\times}10^4$, 32B. R5: $h{=}1$K, $s{=}0.98$, $n{=}10^5$, 4B. Predictions from the first-order model, to be replaced by measurements (Remark~\ref{rem:empirical}).}
\label{tab:motivation}
\end{table}

The table makes two points. The engine default pays for reactive execution, worst where deep high-selectivity pipelines meet LRU eviction that discards KV the pipeline is about to reuse. And planning alone is not enough: breadth-discard is fully planned yet loses up to $2.9\times$ because it forfeits reuse. What captures the gains is planned reuse, and cohorting captures essentially all of it available in this model, provably. Four further fixed shapes (survivor retention, spill-and-return, pipelined mixing, suffix fusion) fall between the rows of Table~\ref{tab:motivation}; Appendix~\ref{app:schedules} gives the full comparison and the closed forms.

\begin{proposition}[Cohorting is near-optimal in the first-order model]
\label{prop:cohort}
Let $\eps$ be the per-token read-to-prefill ratio: the measured time for a cached pass to read one resident token divided by the measured time to prefill one token on the deployed configuration. If every document fits in free HBM, the corpus is cold (no document's KV exists in any tier when execution begins) and available upfront, the device is solely owned, and step cost is additive across requests with these per-token rates, then depth-first cohorting achieves makespan at most $(1 + z\eps + \delta)\,\mathrm{OPT}$, where $\delta$ collects per-step overheads.
\end{proposition}

The intuition is concrete in Example~\ref{ex:contracts}. Every schedule must acquire every contract's KV in HBM at least once before its first evaluation, and with a cold corpus the only source is prefill, so $\mathrm{OPT}$ is lower-bounded by one corpus prefill; a 4K contract costs about 46\,ms of prefill compute at 4B scale (measured: $9.2\,\mu$s/token linear plus the quadratic attention term). Cohorting pays exactly that floor plus $z$ cached passes, and one cached pass reads the contract's 288\,MiB of KV in about 0.41\,ms, so $\eps \approx 6.2 \times 10^{-3}$ and three filters add under two percent. The datasheet approximation $\eps \approx (m/\beta_{\hbm})/(2P_{\Theta}/\mathrm{FLOPS})$ gives $1.5 \times 10^{-3}$ and understates the measured value by $4\times$: the paged-attention read achieves roughly a fifth of spec bandwidth at filter-width suffixes, and its per-token price grows with suffix width (82--141\,ns at $c = 16$--$64$), so $\eps$ is stated at the $c{=}32$ reference. The approximation's structure still argues that $\eps$ improves with scale, since KV bytes per token grow sublinearly in parameters while prefill FLOPs grow linearly, and that it is invariant under tensor sharding, since at degree $\nu$ per-device KV per token is $m/\nu$ while per-device compute per token is $2P_{\Theta}/(\nu\,\mathrm{FLOPS})$ and $\nu$ cancels; both claims are about the approximation and are untested on measured constants, which exist for one unsharded configuration. The single-document hypothesis relaxes by the factor $g$ under sharding.

\begin{corollary}[Warm corpora]
\label{cor:warm}
Suppose instead that when execution begins each document $d$ either has no stored KV or has KV in a tier $k(d)$ with bandwidth $\beta_{k(d)}$, and define its per-token acquisition cost $a_d = \min\!\big(t_{\mathrm{pre}},\, m/\beta_{k(d)}\big)$. Then $\mathrm{OPT} \ge \sum_d h_d\, a_d$, and cohorting that acquires each document once from its cheapest source achieves makespan at most $(1 + z\eps_w + \delta)\,\mathrm{OPT}$, where $\eps_w = t_{\mathrm{read}} / \min_d a_d = \max\big(\eps,\ \max_{k \text{ used}} \beta_k/\beta_{\hbm}\big)$.
\end{corollary}

Warm state is common: it arises whenever several pipelines run over the same corpus and document KV is persisted across queries, which the placement tiers of Section~\ref{sec:placement} make natural. The corollary has two consequences. Fetching from NVMe or host RAM keeps the slack small, $\beta_k/\beta_{\hbm}$ near $2\times 10^{-3}$ and $1.6\times 10^{-2}$ respectively, so cohorting remains near-optimal on corpora warmed from storage, and the floor itself drops, which is the economics of amortizing prefill across queries. But fetching from a peer device over NVLink puts $\beta_k/\beta_{\hbm}$ near $0.3$, and with $z$ cached passes the slack is no longer negligible, so on NVLink-warm corpora scheduling choices matter even when every other hypothesis of Proposition~\ref{prop:cohort} holds.

\begin{remark}[Calibration]
\label{rem:empirical}
Calibration families C1 and C2 (Section~\ref{sec:calib}) estimate the per-token prefill and cached-read times that define $\eps$ on the deployed configuration, and families C4 and C5 test the additivity hypothesis directly. If the calibrated cost model matches measured step latency within relative error $\eta$ over the calibrated range, the bound weakens by a factor of at most $(1+\eta)^2$. Where C4 or C5 reveal a systematic additivity violation, the guarantee does not apply, and the scheduler of Section~\ref{sec:sched} covers those regimes. Calibration has run on the 4B/H100 configuration: $t_{\mathrm{pre}} = 9.2\,\mu$s/token, $t_{\mathrm{read}} = 57$--$101$\,ns/token (suffix-width dependent), $\eps = 6.2\times10^{-3}$, held-out step-model error $\eta = 0.059$ with $96.5\%$ pairwise ranking accuracy; the C4 mixed-step residual is $7.1\%$ and the C5 imbalance residual $1.3\%$, so no systematic additivity violation appears and the envelope stands. The per-query software residue of the planned executor is 26\,ms, measured host-controlled. The entries of Table~\ref{tab:motivation} remain analytical: the strategy comparison itself has not been re-measured.
\end{remark}

\paragraph{Discussion.} Inside the hypotheses of Proposition~\ref{prop:cohort}, the right response to the $9\times$ engine gap is a cohort policy rather than a solver, and our scheduler recovers cohorting, which is a feasible point of its program. The machinery of Sections~\ref{sec:cost} to \ref{sec:placement} is aimed at the hypotheses' failures, each of which is a real deployment: the corpus streams in, so cohort admission is constrained by arrival and latency; residency is contested by other tenants; a document exceeds free HBM, which occurs past roughly 990K tokens at 4B and 328K at 32B per device; step costs carry second-order composition and imbalance effects that first-order arithmetic cannot see and only calibration can measure; or the corpus is warm from earlier queries, where Corollary~\ref{cor:warm} shows the slack grows with the speed of the source tier. Section~\ref{sec:hardness} records the worst-case complexity and the structure that tames it.

\subsection{Overview of Our Approach}
\label{sec:overview}

Our system has an offline phase and a runtime loop. Offline, we calibrate the step cost model on the deployed device configuration (Section~\ref{sec:calib}) and run a stratified pilot over a small document sample to initialize per-bucket selectivity estimates. At runtime, four components interact. The \emph{cost model} scores candidate steps. The \emph{selectivity estimator} maintains a Beta posterior per (filter, bucket) pair, updated after every step. The \emph{scheduler} first partitions the corpus across replicas, then, per replica, re-solves a rolling-horizon program each cycle, with the solve overlapped with GPU execution of the current step. The \emph{tier manager} executes the resulting placements and tracks residency. Each cycle is: observe state, solve, execute one step, read verdicts, update posteriors. Algorithm~\ref{alg:rh} assembles the loop. The components are modular: the cost model can be swapped for a lookup table, and when the hypotheses of Proposition~\ref{prop:cohort} hold, the entire solver can be replaced by the cohort policy, which the program contains as a feasible point.

\section{Complexity}
\label{sec:hardness}

FPS contains two classical NP-hard problems. With atomic suffixes, unbounded memory, and constant step cost, composing steps under the token budget is makespan minimization on a batch machine with non-identical job sizes, NP-hard by reduction from bin packing and studied by Uzsoy~\cite{uzsoy}. With the evaluation sequence fixed, choosing placements is offline generalized caching with page sizes and fetch costs, NP-hard even in the fault model~\cite{cachinghard}. Both statements hold with all selectivities known and equal to one; Appendix~\ref{app:proofs} spells out the reductions as Propositions~\ref{thm:binpack} and~\ref{thm:caching}.

The hardness relies on heterogeneity. Both reductions need many distinct item types: arbitrary suffix lengths in the first, arbitrary page sizes in the second. A real pipeline has $z$ suffix lengths and $L$ length buckets, so our scheduler plans over groups of interchangeable documents rather than over documents (Section~\ref{sec:sched}), the way a query optimizer plans over histogram buckets rather than over tuples~\cite{histograms}: per-group statistics, a representative length and an estimated pass rate, stand in for per-document values, and the bucket count $L$ trades planning resolution against solver size. Grouping does not make FPS polynomial, and we do not solve it exactly. What it buys is a decision space of $O(zL\tau)$ counts independent of $n$, and a relaxation whose rounding loss is at most one document per group per step, which Lemma~\ref{lem:round} bounds by a vanishing fraction of total work at corpus scale.

\section{Step Cost Model}
\label{sec:cost}

The scheduler ranks candidate steps with a predictor of $\Tcyc$, calibrated from controlled measurements. This section gives the model, the reason its form is what it is, and the calibration design; the identifiability result here determines what any such design can learn.

\subsection{Form and an Identifiability Theorem}

The GPU term is
\begin{equation}
\label{eq:costmodel}
\widehat{T}_{\mathrm{gpu}}(\sigma) = b_0 + f_B\big(B(\sigma)\big) + f_P\big(P(\sigma)\big) + \beta_N N,
\end{equation}
with $f_B, f_P$ monotone nondecreasing one-dimensional fits, $\beta_N \ge 0$, one global intercept $b_0 \ge 0$. Host preparation is $\widehat{T}_{\mathrm{host}} = h_0 + h_N N + h_A A + h_R R$ with $A$ the KV block allocations and $R$ the resident KV blocks: the scheduler rebuilds per-step block tables spanning each request's whole context, and calibration shows a plane without $R$ cannot fit prefill and cached cells at once (its error is $5.3\times$ the mean; with $R$, $0.15$). Tier transfers are priced separately from measured bandwidths, because a PCIe copy and a block-table update differ by orders of magnitude and a lumped coefficient predicts neither.

\paragraph{Device configuration enters through constants.} Calibration runs on the deployed configuration of Section~\ref{sec:devices}, so sharding requires no change to the model form: at degree $\nu$ the effective per-token KV is $m/\nu$ per device, compute rates multiply by $\nu$, aggregate free KV grows accordingly, and all-reduce overhead surfaces as measured inflation of $b_0$ and $f_B$. Nothing downstream of calibration distinguishes sharded from unsharded serving.

A natural extension adds a resident-bytes term $R(\sigma) = m\sum_i (H_i + C_i)$ for memory traffic alongside the compute term $P$. We prove this two-term model cannot be learned from the designs this workload naturally produces.

\begin{proposition}[Non-identifiability at a common suffix length]
\label{prop:ident}
In the linear model $T = \theta_0 + \theta_B B + \theta_P P + \theta_R R + \theta_N N$, over any set of calibration cells in which every request has the same new-token count $c$,
\[
P = \tfrac{c}{m} R + \tfrac{c(1-c)}{2} N,
\]
so the design matrix is rank-deficient and $(\theta_P, \theta_R)$ is not identifiable.
\end{proposition}

\paragraph{Discussion.} Filter workloads hold $C_i$ nearly constant within a step by construction, since suffixes span only 16 to 64 tokens, so the hypothesis of Proposition~\ref{prop:ident} is the natural condition, not a pathology. We therefore carry the single attention term $f_P$, exact for compute and proportional to traffic in the small-suffix regime where each cached token is read once. Reintroducing $R$ requires cells built to violate the hypothesis with wide spread in $c$, long suffixes over short contexts against short suffixes over long contexts at matched $P$; these are calibration family C3 below, and they are optional.

\subsection{Calibration Design}
\label{sec:calib}

Every cell submits exactly one step of work with the engine's own scheduler pinned, so measured shapes equal requested shapes, runs five randomized repeats after warmup, and fits the median latency; the p95 is a stability check only. Validation reports mean absolute percentage error and pairwise ranking accuracy over candidate step pairs, by regime. Calibration must boot the engine synchronous and eager: overlapped scheduling moves the GPU wait outside the timed window, and a single large CUDA graph pads small steps up to its capture size. The calibrated $b_0$ therefore carries the eager engine's kernel-launch floor (16.8\,ms at 4B); the production intercept is measured separately with graphs on (2.9\,ms), and the size-dependent terms transfer.

Feasibility bounds the design: the request count at context $h$ is capped by free KV over $h\,m$, about 60 requests at 16K context for 4B and a third of that for 32B. Infeasible cells are excluded upfront, and the model is trusted only inside the measured envelope. $B$ is a derived column, never a swept factor, since Equation~\eqref{eq:BP} makes it a function of composition; a design claiming to sweep $B$ at fixed $(N, \{C_i\})$ is internally inconsistent. Table~\ref{tab:calib} lists the cell families. C4 and C5 exist to measure the second-order effects that Proposition~\ref{prop:cohort} assumes away; if their residuals are systematic, correction terms enter the model and the envelope of the proposition shrinks accordingly.

\begin{table}[t]
\centering
\small
\begin{tabular}{lll}
\toprule
Family & Factors & Identifies \\
\midrule
C1 prefill & $c \in \{64,256,512\}$, $N$ to 512, $H{=}0$ & $f_B$, $\beta_N$, low-$H$ $f_P$ \\
C2 cached & $c \in \{16,32,64\}$, $h$ to 16K & high-$H$ $f_P$ \\
C3 decorr. & matched $P$, varied $R$ & optional $f_R$ \\
C4 mixed & long prefill $+$ cached & mix residual \\
C5 spread & fixed $B,P,N$; varied $H_i$ & imbalance residual \\
C6 transport & PCIe, NVMe, prefill rates & fetch costs (\S\ref{sec:placement}) \\
\bottomrule
\end{tabular}
\caption{Calibration cell families, all within the residency cap.}
\label{tab:calib}
\end{table}

\section{Scheduler}
\label{sec:sched}

The scheduler makes two decisions each cycle: what to run in the next step, and which KV to move between tiers while the step runs. A schedule is this pair repeated (Section~\ref{sec:problem}), and the rest of this section is machinery for choosing the pair well. Choosing it greedily fails for two reasons, both about the future. First, transfers need staging: returning fifty 4K-token survivors from host RAM moves about 14\,GiB, a third of a second at PCIe bandwidth and several steps of overlap, so the KV for a filter that runs at step $t{+}k$ must start moving at step $t$, or the GPU waits. This is the prefetching problem of a buffer manager, with one difference in our favor: the pipeline announces its future reference pattern, since survivors of filter $o$ return at filter $o{+}1$, so eviction and prefetch can approximate Belady's offline rule~\cite{belady} using estimated selectivities, where a request-stream engine can only guess from recency. Second, a one-step objective starves admission: every individual step finishes sooner without a fresh long prefill in it, so a greedy scheduler defers the expensive work indefinitely; the terminal cost in Equation~\eqref{eq:objective} charges deferred work at calibrated throughput so the lookahead cannot hide it. The program below therefore plans $\tau$ steps ahead on expected counts, executes only the first step's decisions, and re-solves; the variables for later steps are a forecast used to price today's choices, discarded after each cycle.

\subsection{Replica Balancing and Grouped State}

The scheduler consumes the device configuration first. Documents are independent and expected per-document work is known per group from bucket length and estimated selectivities, so balancing across the $G/\nu$ replicas is not the hard part: assigning each group's documents round-robin across replicas gives every replica the same count per group up to one document, so replica loads differ by at most $\sum_g w_g$, a relative imbalance of at most $(G/\nu)/\min_g n_g$ by the same argument as Lemma~\ref{lem:round}, negligible at corpus scale. Each replica then runs the loop of this section independently, makespan is the maximum over replicas, and everything that follows is per replica, including Proposition~\ref{prop:cohort}.

Within a replica, documents awaiting the same filter in the same bucket are interchangeable, so state aggregates by group $g = (o,l)$: counts $n_{g,t}$ and tier occupancies $m_{g,k,t}$.

\subsection{The Horizon Program}

Per horizon step $u$, the decisions are evaluations $x_{g,u} \le m_{g,\hbm,u}$, prefill token allocations $q_{l,u}$, and transfer counts $w_{g,k \to k',u}$; they induce the step statistics
\begin{align}
N_u &= \textstyle\sum_g x_{g,u} + r_u, \nonumber\\
B_u &= \textstyle\sum_g x_{g,u} c_{o(g)} + \textstyle\sum_l q_{l,u}, \label{eq:induced}\\
P_u &= \textstyle\sum_g x_{g,u}\big( c_o h_l + \tfrac{c_o(c_o+1)}{2}\big) + P^{\mathrm{pre}}_u, \nonumber
\end{align}
where $r_u$ counts the documents receiving prefill tokens in step $u$ and $P^{\mathrm{pre}}_u$ sums their attention pairs at their positions.
Group dynamics under estimated selectivities are conservation on both sides:
\begin{align}
n_{o,l,u+1} &= n_{o,l,u} - x_{o,l,u}, \label{eq:dyn}\\
n_{o+1,l,u+1} &= n_{o+1,l,u} + \shat_{o,l}\, x_{o,l,u}, \nonumber
\end{align}
with failed documents leaving and freeing KV, and occupancies following the same flows. Constraints per step: the token budget $B_u \le \Bmax$; the HBM budget counting growth,
\begin{equation}
\label{eq:hbm}
\textstyle\sum_g m_{g,\hbm,u} h_l m + m B^{\hbm}_u + W \le M_{\hbm},
\end{equation}
where $B^{\hbm}_u$ is new tokens written in HBM during the step and $W$ is workspace, since budgeting only start-of-step residency admits steps that overflow mid-execution; capacities of the other tiers; and transferred bytes at most bandwidth times cycle time. The objective over a horizon of $\tau$ steps is
\begin{equation}
\label{eq:objective}
\min \sum_{u=t}^{t+\tau-1} \widehat{T}_{\mathrm{cyc}}(\sigma_u) + \Phi(n_{\cdot,t+\tau}),
\quad \Phi(n) = \frac{P^{\mathrm{rem}}(n)}{\rho^{\ast}},
\end{equation}
a terminal cost that charges deferred work at a calibrated pair throughput $\rho^{\ast}$; without it, short horizons push expensive prefill past the horizon edge. When the fitted curves are convex, the relaxation is a convex program; integrality is enforced only on the executed step by rounding fractional counts down, and the loss is provably small.

\begin{lemma}[Rounding loss]
\label{lem:round}
Rounding each step's fractional counts down defers less than one document per group per step, and the deferred work re-enters the next solve. Under additive step cost, the resulting makespan exceeds the fractional plan's by a relative factor of at most $\tau / \min_g n_g$, where $n_g$ is the initial size of group $g$.
\end{lemma}

For the workloads of Section~\ref{sec:eval}, with $\tau = 10$ and thousands of documents per group, the loss is around one percent or below.

\subsection{Solver Overhead and Overlap}
\label{sec:overhead}

The program has $O(zL\tau)$ variables independent of $n$, on the order of $10^3$ for realistic pipelines, solving in single-digit milliseconds cold and well under a millisecond warm-started. Equation~\eqref{eq:cycle} contains the slot that hides this: the solve for step $t{+}1$ runs on the host during GPU execution of step $t$, adding nothing whenever $T_{\mathrm{solve}} + \Thost \le \Tgpu(\sigma_t)$. Prefill-heavy steps of hundreds of milliseconds hide it trivially; the binding case is a cache-heavy step of tens of milliseconds, which a warm-started re-solve fits. Planning one step ahead of observations is the same certainty-equivalent approximation the program makes through $\shat$, and the gap between expected and realized counts is binomial noise absorbed at the next re-solve. Two fallbacks: if an integer solve overruns, an asynchronous full solve every few hundred milliseconds combines with a per-step repair that rescales planned counts by realized verdicts; in the drain phase, when counts are small, a greedy rule (evaluate everything resident, fill with prefill) replaces the solver. The unhidden serial strip, verdict readout and count updates, is microseconds and lives in $\Tser$.

\subsection{Selectivity Estimation}
\label{sec:selest}

Selectivity plays no role inside the envelope of Proposition~\ref{prop:cohort}, where cohorting ignores it. It matters in three places: replica balancing, since expected per-document work depends on how far a document travels through the pipeline; retention and spill decisions, since the horizon program weights a survivor's KV by its probability of reuse through the $\shat$-scaled tier flows of Section~\ref{sec:sched}; and the boundary regimes of Section~\ref{sec:motivation}, where streaming admission and contended placement plan against predicted downstream load. The estimates themselves cost nothing extra: each $(o,l)$ carries a Beta posterior updated from observed verdicts, and the program consumes posterior means. When the corpus is available upfront, a stratified pilot sample estimates all $\shat_{o,l}$ before planning begins, and online updates refine them. Indexing by bucket matters because pass rates correlate with length, and a global estimate biases exactly the downstream counts the horizon plans around. Algorithm~\ref{alg:rh} assembles the loop.

\begin{algorithm}[t]
\caption{Rolling-horizon FPS}
\label{alg:rh}
\KwIn{corpus, device configuration, cost model, horizon $\tau$}
optionally run a stratified pilot to initialize $\shat$\;
partition the corpus across replicas by LPT on expected work; on each replica:\;
\While{$\sum_g n_{g,t} > 0$}{
  update $\shat_{o,l}$ from observed verdicts\;
  solve \eqref{eq:induced}--\eqref{eq:objective} for steps $t,\dots,t{+}\tau{-}1$ (overlapped with step $t{-}1$)\;
  round and execute the first planned step and its placements\;
  observe verdicts; free KV of failures; $t \leftarrow t{+}1$\;
}
\end{algorithm}

\section{Length-Aware Placement}
\label{sec:placement}

Placement is buffer management in which fetch cost depends on where a document's KV lives and on its length:
\begin{equation}
\label{eq:fetch}
T^{\mathrm{fetch}}_k(h) = \frac{h\,m}{\beta_k}, \qquad
T^{\mathrm{fetch}}_{\none}(h) = \widehat{T}_{\mathrm{pre}}(h) = \alpha_1 h + \alpha_2 h^2,
\end{equation}
with $\alpha_1$ the linear MLP work and $\alpha_2 > 0$ the quadratic attention work of recomputation.

\begin{proposition}[Offload crossover]
\label{prop:crossover}
If $\alpha_1 \ge m/\beta_k$, offload to tier $k$ weakly dominates recomputation at every length; otherwise recomputation is cheaper below and offload above the unique threshold $h^{\ast}_k = (m/\beta_k - \alpha_1)/\alpha_2$.
\end{proposition}

\paragraph{Discussion.} The policy this induces is length-aware by derivation: spill long documents to slower tiers, discard and rebuild short ones, and evict shortest-first toward $\none$. The horizon program needs no such rule, since Equation~\eqref{eq:fetch} enters its costs and the behavior emerges, but the threshold explains solutions and serves as the heuristic when running without the solver. Measured at 4B on H100, the first branch of Proposition~\ref{prop:crossover} is the operative one for host memory: $\alpha_1 = 9.2\,\mu$s/token exceeds $m/\beta$ for pinned PCIe ($1.3\,\mu$s at 55\,GB/s) and even unpinned PCIe ($6.8\,\mu$s), so offload to host RAM dominates recomputation at every length; the measured disk tier (3.9\,GB/s read) crosses at about 19{,}900 tokens, past the 16K calibration envelope, so recomputation wins on disk within it. Tier membership follows the same comparison at the model level: for 4B, NVMe beats recomputation only narrowly and $\{\hbm, \cpu\}$ with recomputation as backstop suffices; for 32B, per-token KV roughly doubles, prefill throughput drops by the compute ratio, recomputing a 2K document costs hundreds of milliseconds against tens from NVMe, and host RAM cannot hold a working set, so NVMe becomes the capacity tier. The rule: enable a tier when calibration family C6 shows transfer beating recomputation by a wide margin for the buckets that dominate the corpus and faster tiers lack capacity. Peer-device memory follows the same comparison: when $G/\nu > 1$, a peer replica's free HBM is a tier at NVLink bandwidth, for which the crossover $h^{\ast}$ of Proposition~\ref{prop:crossover} is near zero, so peer memory dominates recomputation at essentially every length, slots above $\cpu$ in the eviction order, and turns a drained replica's HBM into the fastest spill target in the system.

\section{Evaluation Plan}
\label{sec:eval}

\emph{This section states the planned evaluation; all numbers above are analytical placeholders pending these measurements.}

\paragraph{Setup.} One to eight H100s at FP8 KV, on a serving engine modified to accept externally composed steps. Both models fit on a single device, so the parallelism mode is a configuration under test rather than a constraint: Qwen3 4B runs as per-GPU replicas, and Qwen3 32B runs both as per-GPU replicas and as sharded groups of degree 2 and 4, letting us measure the capacity-communication trade-off of Section~\ref{sec:devices} directly. Calibration per Table~\ref{tab:calib} runs once per configuration.

\paragraph{Workloads and the boundary.} The evaluation tests both sides of Proposition~\ref{prop:cohort}. Inside the envelope: batch corpora with heterogeneous lengths and three-to-five filter pipelines, where the prediction is that the solver ties cohorting and both beat the engine default by the margins of Section~\ref{sec:motivation}. Outside: (1) streaming arrival with deadline constraints on admission; (2) a co-tenant occupying a varying fraction of HBM; (3) a long-document corpus at 32B where single documents approach per-device residency; (4) whatever second-order effects C4 and C5 measure, injected back into the analysis.

\paragraph{Baselines.} Engine default; cohorting as the strongest fixed policy; the solver without selectivity estimates; the solver restricted to $\{\hbm\}$; an oracle with true selectivities.

\paragraph{Metrics and ablations.} Makespan; cost model MAPE and ranking accuracy by regime; time in avoidable recomputation; solver overhead per step against the overlap condition of Section~\ref{sec:overhead}; tier sets at both scales against Proposition~\ref{prop:crossover}; per-bucket against global selectivity; terminal cost on and off; data-parallel scaling from one to eight replicas against the balance slack of Section~\ref{sec:sched}.

\section{Related Work}
\label{sec:related}

\paragraph{LLM serving.} Orca introduced iteration-level scheduling~\cite{orca}, vLLM paged KV management~\cite{vllm}, and Sarathi-Serve chunked prefill for throughput-latency balance~\cite{sarathi}. All schedule online streams without pipeline structure, selectivity, or planned placement, the three columns of Table~\ref{tab:approaches}.

\paragraph{Prefix caching and offload.} RadixAttention~\cite{sglang}, CachedAttention~\cite{cachedattention}, and Mooncake~\cite{mooncake} retain and reuse KV with recency policies; FlexGen~\cite{flexgen} plans offload for a fixed request order on one GPU. Our placement is chosen jointly with composition, with fetch costs that separate transfer from recomputation per Proposition~\ref{prop:crossover}.

\paragraph{Semantic operators and cascades.} DocETL, LOTUS, and Palimpzest optimize operator pipelines logically~\cite{docetl,lotus,palimpzest}; SUPG and BARGAIN bound which model answers each record~\cite{supg,bargain}; FrugalGPT selects providers~\cite{frugalgpt}. These reduce the invocation set; FPS executes a fixed set and composes with all of them. Filter reordering by cost and selectivity~\cite{predmig,probpred} is complementary and out of scope.

\paragraph{Query execution models.} The strategy space of Section~\ref{sec:motivation} mirrors classical execution models: operator-at-a-time materialization~\cite{graefe-survey}, block-pipelined execution~\cite{monetx100}, and operator fusion~\cite{hyper}. Our contribution is the analysis of their KV analogues, and the finding that the block-pipelined point is near-optimal for zero-decode pipelines in the first-order model.

\paragraph{Scheduling theory.} Section~\ref{sec:hardness} inherits hardness from batch-machine makespan scheduling~\cite{uzsoy} and generalized caching~\cite{cachinghard}.

\section{Conclusion}
\label{sec:conclusion}

We formalized the batching and placement problem behind zero-decode filter pipelines on one or more GPUs, proved it hard, and then bounded how much the hardness matters: inside a first-order envelope a cohort policy is provably near-optimal while reactive engines lose up to $9\times$, and the paper's machinery, a calibrated cost model with identifiability guarantees, a hidden-latency grouped solver, and a derived length-aware placement policy, targets the envelope's boundary. Open problems include filter reordering, speculative evaluation under high estimated selectivity, and pipeline-parallel sharding, where the per-device KV structure differs from both modes analyzed here.

\appendix
\section{Full Analytical Schedule Comparison}
\label{app:schedules}

Table~\ref{tab:allschedules} extends Table~\ref{tab:motivation} to all fixed shapes of Section~\ref{sec:motivation}. \emph{Breadth-retain} keeps survivor KV resident in HBM and is feasible only when it fits. \emph{Pipelined mix} overlaps the next cohort's prefill with the current cohort's cached evaluations. \emph{Spill-return} pushes survivor KV to lower tiers and streams it back per filter. \emph{Fused} evaluates consecutive suffixes in one hot pass and is sensible only near unit selectivity. The recompute-free, capacity-respecting shapes tie within single-digit percentages, which is the content of Proposition~\ref{prop:cohort}; the shapes that recompute (breadth-discard) or transfer (spill-return) separate by the corresponding fetch costs.

\begin{table}[h]
\centering
\small
\begin{tabular}{l@{\hspace{6pt}}ccccc}
\toprule
Strategy & R1 & R2 & R3 & R4 & R5 \\
\midrule
Depth-first cohorting & 1.00 & 1.00 & 1.00 & 1.00 & 1.00 \\
Pipelined mix     & 1.00 & 1.00 & 1.00 &  --  & 1.00 \\
Fused             &  --  &  --  & 1.00 &  --  & 1.00 \\
Breadth-retain    &  --  & 1.00 &  --  &  --  &  --  \\
Spill-return      & 1.03 & 1.07 & 2.02 & 1.09 & 2.16 \\
Breadth-discard   & 1.05 & 1.75 & 2.71 & 1.75 & 2.94 \\
Engine default    & 1.21 & 1.05 & 3.12 & 2.01 & 3.38 \\
\bottomrule
\end{tabular}
\caption{Analytical slowdown vs.\ best per regime; regimes as in Table~\ref{tab:motivation}. Dashes are infeasible or inapplicable cells.}
\label{tab:allschedules}
\end{table}

\section{Proofs}
\label{app:proofs}

\begin{proposition}[Composition]
\label{thm:binpack}
FPS is NP-hard even with a single memory tier of unbounded capacity, all KV resident, and constant step cost.
\end{proposition}

\paragraph{Proof of Proposition~\ref{thm:binpack}.}
Reduce from \textsc{Bin Packing} with items $a_1,\dots,a_z$, capacity $V$, target $K$. Build an FPS instance with $z$ filters, filter $o$ having suffix length $c_o = a_o$, and $z$ documents where document $d_o$ has passed filters $1$ through $o{-}1$ and awaits filter $o$; selectivities are 1, all KV resident, no prefill. Set $\Bmax = V$ and step cost $b_0 > 0$ constant. Suffix atomicity makes each pending evaluation consume $a_o$ budget atomically, and evaluations of distinct documents are independent, so feasible schedules are assignments of items to capacity-$V$ steps and makespan is $b_0$ times the step count. Makespan at most $b_0 K$ holds if and only if the items pack into $K$ bins. \qed

\begin{proposition}[Placement]
\label{thm:caching}
FPS is NP-hard even when the sequence of filter evaluations is fixed, so that only KV placement decisions remain, and with two tiers.
\end{proposition}

\paragraph{Proof of Proposition~\ref{thm:caching}.}
With the evaluation sequence fixed, each document is a page of $h_l m$ bytes that must be $\hbm$-resident at its evaluation times; serving a nonresident document costs its fetch time. This is offline generalized caching with page sizes and fetch costs, NP-hard even in the fault model~\cite{cachinghard}; the request sequence is the evaluation sequence. \qed

\paragraph{Proof of Proposition~\ref{prop:ident}.}
With $C_i = c$ for all requests in a cell of size $N$ and contexts $H_1,\dots,H_N$: $P = c\sum_i H_i + N\frac{c(c+1)}{2}$ and $R = m(\sum_i H_i + Nc)$. Substituting $\sum_i H_i = R/m - Nc$ gives $P = \frac{c}{m}R + \frac{c(1-c)}{2}N$, an exact linear dependence among the columns for $P$, $R$, and $N$. \qed

\paragraph{Proof of Proposition~\ref{prop:cohort}.}
Write $t_{\mathrm{pre}}$ and $t_{\mathrm{read}}$ for the measured per-token prefill and cached-read times, so $\eps = t_{\mathrm{read}}/t_{\mathrm{pre}}$. Under additive step cost, every feasible schedule evaluates $F_1$ on every document, which requires acquiring that document's KV in HBM; with a cold corpus the only source is prefill, so $\mathrm{OPT} \ge \sum_l n_l h_l t_{\mathrm{pre}}$. Cohorting prefills each document once, transfers nothing, recomputes nothing, and adds at most $z$ cached passes per document; a pass over $h$ resident tokens costs $h\, t_{\mathrm{read}} = \eps\, h\, t_{\mathrm{pre}}$, which is $\eps$ times that document's prefill cost. Summing gives the factor $1 + z\eps$; step overheads contribute $\delta$, at most $b_0$ times the step count over $\mathrm{OPT}$. Feasibility of every cohort uses the single-document hypothesis. \qed

\paragraph{Proof of Corollary~\ref{cor:warm}.}
Acquisition of document $d$'s KV in HBM before its first evaluation costs at least $h_d\, a_d$ under additive cost, whichever source a schedule uses, giving the floor. Cohorting acquires each document exactly once from its cheapest source, paying $\sum_d h_d a_d$, plus at most $z$ cached passes of cost $h_d\, t_{\mathrm{read}} \le \eps_w\, h_d\, a_d$ each, plus step overheads $\delta$. \qed

\paragraph{Proof of Proposition~\ref{prop:crossover}.}
$T^{\mathrm{fetch}}_k(h) - T^{\mathrm{fetch}}_{\none}(h) = (m/\beta_k - \alpha_1)h - \alpha_2 h^2$ is a downward parabola through the origin with positive root $h^{\ast}_k$; positive on $(0, h^{\ast}_k)$, negative beyond. \qed

\paragraph{Proof of Lemma~\ref{lem:round}.}
Rounding down changes each count by less than one, so each step defers less than one document's evaluations or prefill per group, adding less than $\sum_g w_g$ of work per step, where $w_g$ is the per-document work of group $g$; over $\tau$ steps the addition is below $\tau \sum_g w_g$. Total work is at least $\sum_g n_g w_g \ge (\min_g n_g) \sum_g w_g$, and under additive cost makespan is proportional to work up to the overhead term, giving the ratio. \qed

\begin{thebibliography}{10}
\setlength{\itemsep}{1pt}

\bibitem{docetl}
S.~Shankar, T.~Chambers, T.~Shah, A.~G. Parameswaran, and E.~Wu.
\newblock DocETL: Agentic query rewriting and evaluation for complex document processing.
\newblock arXiv:2410.12189, 2024.

\bibitem{lotus}
L.~Patel, S.~Jha, C.~Guestrin, and M.~Zaharia.
\newblock LOTUS: Enabling semantic queries with LLMs over tables of unstructured and structured data.
\newblock arXiv:2407.11418, 2024.

\bibitem{palimpzest}
C.~Liu, M.~Russo, M.~Cafarella, L.~Cao, P.~B. Chen, Z.~Chen, M.~Franklin, T.~Kraska, S.~Madden, and G.~Vitagliano.
\newblock A declarative system for optimizing AI workloads.
\newblock arXiv:2405.14696, 2024.

\bibitem{supg}
D.~Kang, E.~Gan, P.~Bailis, T.~Hashimoto, and M.~Zaharia.
\newblock Approximate selection with guarantees using proxies.
\newblock {\em PVLDB}, 13(11), 2020.

\bibitem{bargain}
S.~Zeighami, S.~Shankar, and A.~Parameswaran.
\newblock Cut costs, not accuracy: LLM-powered data processing with guarantees.
\newblock arXiv:2509.02896, 2025.

\bibitem{taskcascades}
S.~Shankar, S.~Zeighami, and A.~G. Parameswaran.
\newblock Task cascades for efficient unstructured data processing.
\newblock {\em Proc. ACM Manag. Data (SIGMOD)}, 4(1), 2026.

\bibitem{frugalgpt}
L.~Chen, M.~Zaharia, and J.~Zou.
\newblock FrugalGPT: How to use large language models while reducing cost and improving performance.
\newblock arXiv:2305.05176, 2023.

\bibitem{predmig}
J.~M. Hellerstein and M.~Stonebraker.
\newblock Predicate migration: Optimizing queries with expensive predicates.
\newblock In {\em SIGMOD}, 1993.

\bibitem{probpred}
Y.~Lu, A.~Chowdhery, S.~Kandula, and S.~Chaudhuri.
\newblock Accelerating machine learning inference with probabilistic predicates.
\newblock In {\em SIGMOD}, 2018.

\bibitem{vllm}
W.~Kwon, Z.~Li, S.~Zhuang, Y.~Sheng, L.~Zheng, C.~H. Yu, J.~Gonzalez, H.~Zhang, and I.~Stoica.
\newblock Efficient memory management for large language model serving with PagedAttention.
\newblock In {\em SOSP}, 2023.

\bibitem{orca}
G.-I. Yu, J.~S. Jeong, G.-W. Kim, S.~Kim, and B.-G. Chun.
\newblock Orca: A distributed serving system for transformer-based generative models.
\newblock In {\em OSDI}, 2022.

\bibitem{sarathi}
A.~Agrawal, N.~Kedia, A.~Panwar, J.~Mohan, N.~Kwatra, B.~Gulavani, A.~Tumanov, and R.~Ramjee.
\newblock Taming throughput-latency tradeoff in LLM inference with Sarathi-Serve.
\newblock In {\em OSDI}, 2024.

\bibitem{sglang}
L.~Zheng, L.~Yin, Z.~Xie, C.~Sun, J.~Huang, C.~H. Yu, S.~Cao, C.~Kozyrakis, I.~Stoica, J.~E. Gonzalez, C.~Barrett, and Y.~Sheng.
\newblock SGLang: Efficient execution of structured language model programs.
\newblock In {\em NeurIPS}, 2024.

\bibitem{cachedattention}
B.~Gao, Z.~He, P.~Sharma, Q.~Kang, D.~Jevdjic, J.~Deng, X.~Yang, Z.~Yu, and P.~Zuo.
\newblock Cost-efficient large language model serving for multi-turn conversations with CachedAttention.
\newblock In {\em USENIX ATC}, 2024.

\bibitem{mooncake}
R.~Qin, Z.~Li, W.~He, M.~Zhang, Y.~Wu, W.~Zheng, and X.~Xu.
\newblock Mooncake: A KVCache-centric disaggregated architecture for LLM serving.
\newblock In {\em FAST}, 2025.

\bibitem{flexgen}
Y.~Sheng, L.~Zheng, B.~Yuan, Z.~Li, M.~Ryabinin, B.~Chen, P.~Liang, C.~R\'e, I.~Stoica, and C.~Zhang.
\newblock FlexGen: High-throughput generative inference of large language models with a single GPU.
\newblock In {\em ICML}, 2023.

\bibitem{graefe-survey}
G.~Graefe.
\newblock Query evaluation techniques for large databases.
\newblock {\em ACM Computing Surveys}, 25(2):73--169, 1993.

\bibitem{monetx100}
P.~A. Boncz, M.~Zukowski, and N.~Nes.
\newblock MonetDB/X100: Hyper-pipelining query execution.
\newblock In {\em CIDR}, 2005.

\bibitem{hyper}
T.~Neumann.
\newblock Efficiently compiling efficient query plans for modern hardware.
\newblock {\em PVLDB}, 4(9), 2011.

\bibitem{belady}
L.~A. Belady.
\newblock A study of replacement algorithms for a virtual-storage computer.
\newblock {\em IBM Systems Journal}, 5(2):78--101, 1966.

\bibitem{histograms}
Y.~Ioannidis.
\newblock The history of histograms (abridged).
\newblock In {\em VLDB}, 2003.

\bibitem{uzsoy}
R.~Uzsoy.
\newblock Scheduling a single batch processing machine with non-identical job sizes.
\newblock {\em International Journal of Production Research}, 32(7):1615--1635, 1994.

\bibitem{cachinghard}
M.~Chrobak, G.~J. Woeginger, K.~Makino, and H.~Xu.
\newblock Caching is hard, even in the fault model.
\newblock In {\em ESA}, 2010.

\end{thebibliography}

\end{document}
